
\chapter*{ 致谢}
\addcontentsline{toc}{chapter}{致谢}
写作是一项孤独的事业。形单影只，与键盘为伴，上一秒文思泉涌，下一秒却自我怀疑。

但从另一些方面来说，写作又是一种高度社会化的工作。每一位作者都可以回头追问：是谁点燃了我们的创作欲望？是谁在关键时刻给出了那一句提醒？我们又曾从哪些作家、思想家、研究者以及各种不同的思想中汲取了营养。正是这些可见与不可见的同行者，让一本书慢慢长成现在的样子。

下列个人与机构都给予了我们特别的启发与帮助，谨致谢忱。
 
%我们无比感谢本书的中英文版稿件贡献者和论坛用户们。他们帮助增添或改进了书中内容并提供了有价值的反馈。特别地，我们要感谢每一位为这本中文版开源书提交内容改动的贡献者们。这些贡献者的GitHub用户名或姓名是（排名不分先后）：.....谢谢你们为每一位读者改进这本开源书。
本书的初稿在华东师范大学软件工程学院的“人工智能的数学思维”、华东师范大学通视课程《》中被用于教学。
我们在此感谢这些课程的师生，特别是  ，感谢他们对改进本书提供的宝贵意见。
此外，我们感谢
%Amazon Web Services，特别是Swami Sivasubramanian、Raju Gulabani、Charlie Bell和Andrew Jassy
在我们撰写本书时给予的慷慨支持。如果没有可用的时间、资源以及来自同事们的讨论和鼓励，就没有这本书的项目。我们还要感谢Apache MXNet团队实现了很多本书所使用的特性。

另外，经过同事们的校勘，本书的质量得到了极大的提升。在此我们一一列出章节和校勘人，以表示我们由衷的感谢：引言的校勘人为金颢，预备知识的校勘人为吴俊，深度学习基础的校勘人为张航、王晨光、林海滨，深度学习计算的校勘人为查晟，卷积神经网络的校勘人为张帜、何通，循环神经网络的校勘人为查晟，优化算法的校勘人为郑帅，计算性能的校勘人为郑达、吴俊，计算机视觉的校勘人为解浚源、张帜、何通、张航，自然语言处理的校勘人为王晨光，附录的校勘人为金颢。
%
%
%感谢将门创投，特别是王慧、高欣欣、常铭珊和白玉为本书的两位中国作者在讲授”动手学深度学习“系列课程时所提供的平台。感谢所有参与这一系列课程的数千名同学们。感谢Amazon Web Services中国团队的同事们，特别是费良宏和王晨对作者的支持与鼓励。感谢本书论坛的3位版主：王鑫、夏鲁豫和杨培文。他们牺牲了自己宝贵的休息时间来回复大家的提问。感谢人民邮电出版社的杨海玲编辑为我们在本书的出版过程中所提供的各种帮助。
%感谢中英文草稿的数百位撰稿人。
%他们帮助改进了内容并提供了宝贵的反馈。
%感谢Anirudh Dagar和唐源将部分较早版本的MXNet实现分别改编为PyTorch和TensorFlow实现。
%感谢百度团队将较新的PyTorch实现改编为PaddlePaddle实现。
%感谢张帅将更新的LaTeX样式集成进PDF文件的编译。
%特别地，我们要感谢这份中文稿的每一位撰稿人，是他们的无私奉献让这本书变得更好。他们的GitHub ID或姓名是(没有特定顺序)：
%我们感谢Amazon Web Services，特别是Swami Sivasubramanian、Peter DeSantis、Adam Selipsky和Andrew Jassy对撰写本书的慷慨支持。如果没有可用的时间、资源、与同事的讨论和不断的鼓励，这本书就不会出版。

最后，我要特别感谢我们的家人。谢谢你们一直默默陪伴着我们，在我们埋头写作、反复修改、偶尔沮丧的日子里，承担了许多看不见的时间与情绪成本。没有你们的理解与支持，这本书无法按现在的样子与大家见面。
